\documentclass[12pt,aspectratio=169,ignorenonframetext]{ltxbeamer}

\usepackage[spanish,es-nodecimaldot,es-tabla]{babel}
\usepackage{calc}
\usetikzlibrary{arrows.meta}
\usepackage{hologo}

\title{Matemáticas}
\subtitle{Curso de {\lmr\LaTeX}}
\author{Joar Buitrago\texorpdfstring{\thanks{jebuitragoc@unal.edu.co}}{ <jebuitragoc@unal.edu.co>}}
\date{}

\addbibresource{bibliography.bib}





\newlength{\mathspacetest}

\begin{document}
\begin{frame}
\titlepage
\end{frame}



\begin{frame}
\frametitle{Matemáticas en {\lmr\LaTeX}}

{\lmr\LaTeX} es muy reconocido por su increíble y sencillo manejo de escritura matemática.

\pause\bigskip

\begin{alertblock}{IMPORTANTE}
    Todos los entornos o comandos de matemáticos de {\lmr\LaTeX} hacen que este entre en modo matemático.
\end{alertblock}
\end{frame}





\begin{frame}
\frametitle{Matemáticas en {\lmr\LaTeX}}

El modo matemático hace que {\lmr\LaTeX} utilice la fuente configurada en ese momento como matemática.

\pause

Este modo se comporta de manera similar que el modo ID. Con la excepción de que se eliminan todos los espacios en blanco del entorno.

\pause\bigskip

\begingroup
\scriptsize
\makebox[\linewidth][l]{\(\displaystyle\int\frac{1}{x^n + 1}\,dx = - \frac{1}{n} \sum_{k=1}^n \left\{ \frac{1}{2} \cos \left( \frac{2 \pi k - \pi}{n} \right) \ln \left[ x^2 - 2 \cos \left( \frac{2 \pi k - \pi}{n} \right) x + 1 \right] - \sin \left( \frac{2 \pi k - \pi}{n} \right) \arctan \left[ \frac{ x - \cos \left( \frac{2 \pi k - \pi}{n} \right)}{\sin \left( \frac{2 \pi k - \pi}{n} \right)} \right] \right\} + C\)}
\endgroup
\end{frame}





\section{Primitivas de {\lmr\TeX}}

\begin{frame}[fragile]
\frametitle{Primitivas de {\lmr\TeX}}

{\lmr\TeX} ofrece dos primitivas o comandos para escribir ecuaciones.

\pause\medskip

\begin{columns}
    \begin{column}{0.5\textwidth}
        \begin{minted}{latex}
$a^2 + b^2 = c^2$
        \end{minted}
    \end{column}
    {\color{gray} \vline}
    \begin{column}{0.5\textwidth}
        $a^2 + b^2 = c^2$
    \end{column}
\end{columns}

\pause\medskip

\begin{columns}
    \begin{column}{0.5\textwidth}
        \begin{minted}{latex}
$$a^2 + b^2 = c^2$$
        \end{minted}
    \end{column}
    {\color{gray} \vline}
    \begin{column}{0.5\textwidth}
        $$a^2 + b^2 = c^2$$
    \end{column}
\end{columns}

\pause\bigskip

\begin{alertblock}{IMPORTANTE}
    Son primitivas de {\lmr\TeX}.

    No deberían usarse en un documento de {\lmr\LaTeX}.
\end{alertblock}
\end{frame}





\section{Primitivas de {\lmr\LaTeX}}

\begin{frame}[fragile]
\frametitle{Primitivas de {\lmr\LaTeX}}

Los comandos más básicos que ofrece {\lmr\LaTeX} son igualmente dos.

\pause\medskip

\begin{columns}
    \begin{column}{0.5\textwidth}
        \begin{minted}{latex}
\(a^2 + b^2 = c^2\)
        \end{minted}
    \end{column}
    {\color{gray} \vline}
    \begin{column}{0.5\textwidth}
        \(a^2 + b^2 = c^2\)
    \end{column}
\end{columns}

\pause\medskip

\begin{columns}
    \begin{column}{0.5\textwidth}
        \begin{minted}{latex}
\[a^2 + b^2 = c^2\]
        \end{minted}
    \end{column}
    {\color{gray} \vline}
    \begin{column}{0.5\textwidth}
        \[a^2 + b^2 = c^2\]
    \end{column}
\end{columns}
\end{frame}





\begin{frame}[fragile]
\frametitle{Diferencia entre \texttt{inline} y \texttt{display}}

\pause\medskip

\begin{columns}
    \begin{column}{0.5\textwidth}
        \centering
        \mintinline{latex}{\(\frac{1}{2}\)}\par\bigskip
        \(\frac{1}{2}\)\par
    \end{column}
    \begin{column}{0.5\textwidth}
        \centering
        \mintinline{latex}{\[\frac{1}{2}\]}\par\bigskip
        \(\displaystyle\frac{1}{2}\)\par
    \end{column}
\end{columns}
\end{frame}





\begin{frame}[fragile]
\frametitle{Entornos matemáticos}

{\lmr\LaTeX} nos ofrece dos entornos matemáticos, ambos con el estilo \texttt{display}, en su propia línea.

\pause\medskip

\begin{columns}
    \begin{column}{0.5\textwidth}
        \begin{minted}{latex}
\begin{equation}
    a^2 + b^2 = c^2
\end{equation}
        \end{minted}
    \end{column}
    {\color{gray} \vline}
    \begin{column}{0.5\textwidth}
        \lmr\strut\hfill\(\displaystyle a^2 + b^2 = c^2\)\hfill\makebox[0pt][r]{(1)}
    \end{column}
\end{columns}

\pause\medskip

\begin{columns}
    \begin{column}{0.5\textwidth}
        \begin{minted}{latex}
\begin{displaymath}
    a^2 + b^2 = c^2
\end{displaymath}
        \end{minted}
    \end{column}
    {\color{gray} \vline}
    \begin{column}{0.5\textwidth}
        \begin{displaymath}
            a^2 + b^2 = c^2
        \end{displaymath}
    \end{column}
\end{columns}
\end{frame}





\begin{frame}[fragile]
\frametitle{Entornos matemáticos}

La primitiva \mintinline{latex}{\[\]} en una abreviatura del entorno \texttt{displaymath}.

Por lo que con ambos se obtendrá siempre el mismo resultado.\pause\alert{*}
\end{frame}





\begin{frame}[fragile]
\frametitle{Entornos matemáticos}

El entorno \texttt{equation} nos permite numerar las ecuaciones.
Que podemos etiquetar y referenciar luego con las órdenes \mintinline{latex}{\label} y \mintinline{latex}{\ref}.

\pause\bigskip

\begin{columns}
    \begin{column}{0.5\textwidth}
        \begin{minted}{latex}
\begin{equation}
    a^2 + b^2 = c^2
    \label{eq:pitagoras}
\end{equation}

Como podemos ver en
la ecuación \ref{eq:pitagoras}.
        \end{minted}
    \end{column}
    {\color{gray} \vline}
    \begin{column}{0.5\textwidth}
        \lmr
        \strut\hfill\(\displaystyle a^2 + b^2 = c^2\)\hfill\makebox[0pt][r]{(1)}\par\medskip

        Como podemos ver en la ecuación 1.
    \end{column}
\end{columns}
\end{frame}





\section{Algunos comandos matemáticos}

\begin{frame}[fragile]
\frametitle{Algunos comando matemáticos}

\begin{columns}
    \begin{column}{0.5\textwidth}
        \begin{minted}{latex}
\[\frac{up}{down}\]
        \end{minted}
    \end{column}
    {\color{gray} \vline}
    \begin{column}{0.5\textwidth}
        \[\frac{up}{down}\]
    \end{column}
\end{columns}
\end{frame}





\begin{frame}[fragile]
\frametitle{Algunos comando matemáticos}

\begin{columns}
    \begin{column}{0.5\textwidth}
        \begin{minted}{latex}
\[\frac{\frac{1}{2} + 2}{3}\]
        \end{minted}
    \end{column}
    {\color{gray} \vline}
    \begin{column}{0.5\textwidth}
        \[\frac{\frac{1}{2} + 2}{3}\]
    \end{column}
\end{columns}
\end{frame}





\begin{frame}[fragile]
\frametitle{Algunos comando matemáticos}

\begin{columns}
    \begin{column}{0.5\textwidth}
        \begin{minted}{latex}
\[\left(  \right)\]
        \end{minted}
    \end{column}
    {\color{gray} \vline}
    \begin{column}{0.5\textwidth}
        \[\left(  \right)\]
    \end{column}
\end{columns}
\end{frame}





\begin{frame}[fragile]
\frametitle{Algunos comando matemáticos}

\begin{columns}
    \begin{column}{0.5\textwidth}
        \begin{minted}{latex}
\[\left( \frac{1}{2} \right)\]
        \end{minted}
    \end{column}
    {\color{gray} \vline}
    \begin{column}{0.5\textwidth}
        \[\left( \frac{1}{2} \right)\]
    \end{column}
\end{columns}
\end{frame}





\begin{frame}[fragile]
\frametitle{Algunos comando matemáticos}

\begin{columns}
    \begin{column}{0.5\textwidth}
        \begin{minted}{latex}
\[\left[ \frac{1}{2} \right\}\]
        \end{minted}
    \end{column}
    {\color{gray} \vline}
    \begin{column}{0.5\textwidth}
        \[\left[ \frac{1}{2} \right\}\]
    \end{column}
\end{columns}
\end{frame}





\begin{frame}[fragile]
\frametitle{Algunos comando matemáticos}

\begin{columns}
    \begin{column}{0.5\textwidth}
        \begin{minted}{latex}
\[\left] \frac{1}{2} \right.\]
        \end{minted}
    \end{column}
    {\color{gray} \vline}
    \begin{column}{0.5\textwidth}
        \[\left] \frac{1}{2} \right.\]
    \end{column}
\end{columns}
\end{frame}





\begin{frame}[fragile]
\frametitle{Algunos comando matemáticos}

\begingroup
\small
\centering
\begin{tabular}{lc}
    \toprule
    Símbolo & Visualización \\
    \midrule
    \mintinline{latex}{()} & \(\displaystyle\left(\rule{0pt}{1em}\right)\) \\[1em]
    \mintinline{latex}{[]} o \mintinline{latex}{\lbrack\rbrack} & \(\displaystyle\left[\rule{0pt}{1em}\right]\) \\[1em]
    \mintinline{latex}{\{\}} o \mintinline{latex}{\lbrace\rbrace} & \(\displaystyle\left\{\rule{0pt}{1em}\right\}\) \\[1em]
    \mintinline{latex}{\lfloor\rfloor} & \(\displaystyle\left\lfloor\rule{0pt}{1em}\right\rfloor\) \\[1em]
    \mintinline{latex}{\lceil\rceil} & \(\displaystyle\left\lceil\rule{0pt}{1em}\right\rceil\) \\[1em]
    \mintinline{latex}{\langle\rangle} & \(\displaystyle\left\langle\rule{0pt}{1em}\right\rangle\) \\[1em]
    \mintinline{latex}{\lgroup\rgroup} & \(\displaystyle\left\lgroup\rule{0pt}{1em}\right\rgroup\) \\
    \bottomrule
\end{tabular}
\par
\endgroup
\end{frame}





\begin{frame}[fragile]
\frametitle{Algunos comando matemáticos}

\begin{columns}
    \begin{column}{0.5\textwidth}
        \small
        \centering
        \begin{tabular}{lc}
            \toprule
            Símbolo & Visualización \\
            \midrule
            \mintinline{latex}{/} & \(\displaystyle\left/\rule{0pt}{1em}\right.\) \\[1em]
            \mintinline{latex}{\backslash} & \(\displaystyle\left\backslash\rule{0pt}{1em}\right.\) \\[1em]
            \mintinline{latex}{|} o \mintinline{latex}{\vert} & \(\displaystyle\left|\rule{0pt}{1em}\right.\) \\[1em]
            \mintinline{latex}{\|} o \mintinline{latex}{\Vert} & \(\displaystyle\left\|\rule{0pt}{1em}\right.\) \\
            \bottomrule
        \end{tabular}
        \par
    \end{column}
    \begin{column}{0.5\textwidth}
        \small
        \centering
        \begin{tabular}{lc}
            \toprule
            Símbolo & Visualización \\
            \midrule
            \mintinline{latex}{\uparrow} & \(\displaystyle\left\uparrow\rule{0pt}{1em}\right.\) \\[1em]
            \mintinline{latex}{\Uparrow} & \(\displaystyle\left\Uparrow\rule{0pt}{1em}\right.\) \\[1em]
            \mintinline{latex}{\downarrow} & \(\displaystyle\left\downarrow\rule{0pt}{1em}\right.\) \\[1em]
            \mintinline{latex}{\Downarrow} & \(\displaystyle\left\Downarrow\rule{0pt}{1em}\right.\) \\[1em]
            \mintinline{latex}{\updownarrow} & \(\displaystyle\left\updownarrow\rule{0pt}{1em}\right.\) \\[1em]
            \mintinline{latex}{\Updownarrow} & \(\displaystyle\left\Updownarrow\rule{0pt}{1em}\right.\) \\
            \bottomrule
        \end{tabular}
        \par
    \end{column}
\end{columns}
\end{frame}





\begin{frame}[fragile]
\frametitle{Algunos comando matemáticos}

\begin{columns}
    \begin{column}{0.5\textwidth}
        \begin{minted}{latex}
\[2^2\]
        \end{minted}
    \end{column}
    {\color{gray} \vline}
    \begin{column}{0.5\textwidth}
        \[2^2\]
    \end{column}
\end{columns}
\end{frame}





\begin{frame}[fragile]
\frametitle{Algunos comando matemáticos}

\begin{columns}
    \begin{column}{0.5\textwidth}
        \begin{minted}{latex}
\[2^{algo}\]
        \end{minted}
    \end{column}
    {\color{gray} \vline}
    \begin{column}{0.5\textwidth}
        \[2^{algo}\]
    \end{column}
\end{columns}
\end{frame}





\begin{frame}[fragile]
\frametitle{Algunos comando matemáticos}

\begin{columns}
    \begin{column}{0.5\textwidth}
        \begin{minted}{latex}
\[2^{2^2}\]
        \end{minted}
    \end{column}
    {\color{gray} \vline}
    \begin{column}{0.5\textwidth}
        \[2^{2^2}\]
    \end{column}
\end{columns}
\end{frame}





\begin{frame}[fragile]
\frametitle{Algunos comando matemáticos}

\begin{columns}
    \begin{column}{0.5\textwidth}
        \begin{minted}{latex}
\[2^2^2\]
        \end{minted}
    \end{column}
    {\color{gray} \vline}
    \begin{column}{0.5\textwidth}
        \[2^2{}^2\]
    \end{column}
\end{columns}

\pause

\begin{minted}{text}
! Double superscript.
l.14     2^2^
           2
I treat `x^1^2' essentially like `x^1{}^2'.
\end{minted}
\end{frame}





\begin{frame}[fragile]
\frametitle{Algunos comando matemáticos}

\begin{columns}
    \begin{column}{0.5\textwidth}
        \begin{minted}{latex}
\[2_2\]
        \end{minted}
    \end{column}
    {\color{gray} \vline}
    \begin{column}{0.5\textwidth}
        \[2_2\]
    \end{column}
\end{columns}
\end{frame}





\begin{frame}[fragile]
\frametitle{Algunos comando matemáticos}

\begin{columns}
    \begin{column}{0.5\textwidth}
        \begin{minted}{latex}
\[2_{algo}\]
        \end{minted}
    \end{column}
    {\color{gray} \vline}
    \begin{column}{0.5\textwidth}
        \[2_{algo}\]
    \end{column}
\end{columns}
\end{frame}





\begin{frame}[fragile]
\frametitle{Algunos comando matemáticos}

\begin{columns}
    \begin{column}{0.5\textwidth}
        \begin{minted}{latex}
\[2_{2_2}\]
        \end{minted}
    \end{column}
    {\color{gray} \vline}
    \begin{column}{0.5\textwidth}
        \[2_{2_2}\]
    \end{column}
\end{columns}
\end{frame}





\begin{frame}[fragile]
\frametitle{Algunos comando matemáticos}

\begin{columns}
    \begin{column}{0.5\textwidth}
        \begin{minted}{latex}
\[2_2_2\]
        \end{minted}
    \end{column}
    {\color{gray} \vline}
    \begin{column}{0.5\textwidth}
        \[2_2{}_2\]
    \end{column}
\end{columns}

\pause

\begin{minted}{text}
! Double subscript.
l.14     2_2_
           2
I treat `x_1_2' essentially like `x_1{}_2'.
\end{minted}
\end{frame}





\begin{frame}[fragile]
\frametitle{Algunos comando matemáticos}

\begin{columns}
    \begin{column}{0.5\textwidth}
        \begin{minted}{latex}
\[2 \sp 2 \sb 2\]
        \end{minted}
    \end{column}
    {\color{gray} \vline}
    \begin{column}{0.5\textwidth}
        \[2 \sp 2 \sb 2\]
    \end{column}
\end{columns}
\end{frame}





\begin{frame}[fragile]
\frametitle{Algunos comando matemáticos}

\begin{columns}
    \begin{column}{0.5\textwidth}
        \begin{minted}{latex}
\[\sqrt{x}\]
        \end{minted}
    \end{column}
    {\color{gray} \vline}
    \begin{column}{0.5\textwidth}
        \[\sqrt{x}\]
    \end{column}
\end{columns}
\end{frame}





\begin{frame}[fragile]
\frametitle{Algunos comando matemáticos}

\begin{columns}
    \begin{column}{0.5\textwidth}
        \begin{minted}{latex}
\[\sqrt{\frac{1}{2}}\]
        \end{minted}
    \end{column}
    {\color{gray} \vline}
    \begin{column}{0.5\textwidth}
        \[\sqrt{\frac{1}{2}}\]
    \end{column}
\end{columns}
\end{frame}





\begin{frame}[fragile]
\frametitle{Algunos comando matemáticos}

\begin{columns}
    \begin{column}{0.5\textwidth}
        \begin{minted}{latex}
\[\sqrt{\frac{1}{2} + 2}\]
        \end{minted}
    \end{column}
    {\color{gray} \vline}
    \begin{column}{0.5\textwidth}
        \[\sqrt{\frac{1}{2} + 2}\]
    \end{column}
\end{columns}
\end{frame}





\begin{frame}[fragile]
\frametitle{Algunos comando matemáticos}

\begin{columns}
    \begin{column}{0.5\textwidth}
        \begin{minted}{latex}
\[\sqrt[6]{\frac{1}{2} + 2}\]
        \end{minted}
    \end{column}
    {\color{gray} \vline}
    \begin{column}{0.5\textwidth}
        \[\sqrt[6]{\frac{1}{2} + 2}\]
    \end{column}
\end{columns}
\end{frame}





\begin{frame}[fragile]
\frametitle{Algunos comando matemáticos}

\begin{columns}
    \begin{column}{0.5\textwidth}
        \begin{minted}{latex}
\[\sum i\]
        \end{minted}
    \end{column}
    {\color{gray} \vline}
    \begin{column}{0.5\textwidth}
        \[\sum i\]
    \end{column}
\end{columns}
\end{frame}





\begin{frame}[fragile]
\frametitle{Algunos comando matemáticos}

\begin{columns}
    \begin{column}{0.5\textwidth}
        \begin{minted}{latex}
\[\sum_{i=1}^{n} i\]
        \end{minted}
    \end{column}
    {\color{gray} \vline}
    \begin{column}{0.5\textwidth}
        \[\sum_{i=1}^{n} i\]
    \end{column}
\end{columns}
\end{frame}





\begin{frame}[fragile]
\frametitle{Algunos comando matemáticos}

\begin{columns}
    \begin{column}{0.5\textwidth}
        \begin{minted}{latex}
\[\prod_{i=1}^{n} i\]
        \end{minted}
    \end{column}
    {\color{gray} \vline}
    \begin{column}{0.5\textwidth}
        \[\prod_{i=1}^{n} i\]
    \end{column}
\end{columns}
\end{frame}





\begin{frame}[fragile]
\frametitle{Algunos comando matemáticos}

\begin{columns}
    \begin{column}{0.5\textwidth}
        \begin{minted}{latex}
\[\int_{a}^{b} x \, dx\]
        \end{minted}
    \end{column}
    {\color{gray} \vline}
    \begin{column}{0.5\textwidth}
        \[\int_{a}^{b} x \, dx\]
    \end{column}
\end{columns}
\end{frame}





\begin{frame}[fragile]
\frametitle{Algunos comando matemáticos}

\begin{columns}
    \begin{column}{0.5\textwidth}
        \begin{minted}{latex}
\[\iint_{a}^{b} x \, dx\]
        \end{minted}
    \end{column}
    {\color{gray} \vline}
    \begin{column}{0.5\textwidth}
        \[\iint_{a}^{b} x \, dx\]
        \pause
        \[\int\int_{a}^b x \, dx\]
    \end{column}
\end{columns}
\end{frame}





\begin{frame}[fragile]
\frametitle{Algunos comando matemáticos}

\begin{columns}
    \begin{column}{0.5\textwidth}
        \begin{minted}{latex}
\[\iiint_{a}^{b} x \, dx\]
        \end{minted}
    \end{column}
    {\color{gray} \vline}
    \begin{column}{0.5\textwidth}
        \[\iiint_{a}^{b} x \, dx\]
        \pause
        \[\int\int\int_{a}^b x \, dx\]
    \end{column}
\end{columns}
\end{frame}





\begin{frame}[fragile]
\frametitle{Algunos comando matemáticos}

\begin{columns}
    \begin{column}{0.5\textwidth}
        \begin{minted}{latex}
\[\oint_{a}^{b} x \, dx\]
        \end{minted}
    \end{column}
    {\color{gray} \vline}
    \begin{column}{0.5\textwidth}
        \[\oint_{a}^{b} x \, dx\]
    \end{column}
\end{columns}
\end{frame}





\begin{frame}[fragile]
\frametitle{Espacio Horizontal}

\begingroup
\small
\centering
\begin{tabular}{lcc}
    \toprule
    Símbolo & Tamaño & Visualización \\
    \midrule
    && \(HH\) \\
    \mintinline{latex}{\,} & \settowidth{\mathspacetest}{\(\,\)}\tikz[baseline]{\draw [{|[right,width=2pt]}-{|[left,width=2pt]}] (0,0) -- (\mathspacetest,0);} & \(H\,H\) \\
    \mintinline{latex}{\>} o \mintinline{latex}{\:} & \settowidth{\mathspacetest}{\(\>\)}\tikz[baseline]{\draw [{|[right,width=2pt]}-{|[left,width=2pt]}] (0,0) -- (\mathspacetest,0);} & \(H\>H\) \\
    \mintinline{latex}{\;} & \settowidth{\mathspacetest}{\(\;\)}\tikz[baseline]{\draw [{|[right,width=2pt]}-{|[left,width=2pt]}] (0,0) -- (\mathspacetest,0);} & \(H\;H\) \\
    \mintinline[escapeinside=||,showspaces]{latex}{\ |\strut|} & \settowidth{\mathspacetest}{\(\ \)}\tikz[baseline]{\draw [{|[right,width=2pt]}-{|[left,width=2pt]}] (0,0) -- (\mathspacetest,0);} & \(H\ H\) \\
    \mintinline{latex}{\quad} & \settowidth{\mathspacetest}{\(\quad\)}\tikz[baseline]{\draw [{|[right,width=2pt]}-{|[left,width=2pt]}] (0,0) -- (\mathspacetest,0);} & \(H\quad H\) \\
    \mintinline{latex}{\qquad} & \settowidth{\mathspacetest}{\(\qquad\)}\tikz[baseline]{\draw [{|[right,width=2pt]}-{|[left,width=2pt]}] (0,0) -- (\mathspacetest,0);} & \(H\qquad H\) \\
    \mintinline{latex}{\!} & \settowidth{\mathspacetest}{\(\!\)}\tikz[baseline]{\draw [{|[left,width=2pt]}-{|[right,width=2pt]}] (0,0) -- (\mathspacetest,0);} & \(H\!H\) \\
    \bottomrule
\end{tabular}
\par
\endgroup
\end{frame}





\begin{frame}[fragile]
\frametitle{Espacio Horizontal}

\begingroup
\small
\centering
\begin{tabular}{lcc}
    \toprule
    Símbolo & Tamaño & Visualización \\
    \midrule
    && \(HH\) \\
    \mintinline{latex}{\hspace{2em}} & \settowidth{\mathspacetest}{\(\hspace{2em}\)}\tikz[baseline]{\draw [{|[right,width=2pt]}-{|[left,width=2pt]}] (0,0) -- (\mathspacetest,0);} & \(H\hspace{2em}H\) \\
    \bottomrule
\end{tabular}
\par
\endgroup
\end{frame}





\begin{frame}[fragile]
\frametitle{Operadores}

\begin{columns}
    \begin{column}{0.5\textwidth}
        \begin{minted}{latex}
\[\sin x\]
        \end{minted}
    \end{column}
    {\color{gray} \vline}
    \begin{column}{0.5\textwidth}
        \[\sin x\]
    \end{column}
\end{columns}
\end{frame}





\begin{frame}[fragile]
\frametitle{Operadores}

\begin{columns}
    \begin{column}{0.5\textwidth}
        \begin{minted}{latex}
\[\cos x\]
        \end{minted}
    \end{column}
    {\color{gray} \vline}
    \begin{column}{0.5\textwidth}
        \[\cos x\]
    \end{column}
\end{columns}
\end{frame}





\begin{frame}[fragile]
\frametitle{Operadores}

\begin{columns}
    \begin{column}{0.5\textwidth}
        \begin{minted}{latex}
\[\exp x\]
        \end{minted}
    \end{column}
    {\color{gray} \vline}
    \begin{column}{0.5\textwidth}
        \[\exp x\]
    \end{column}
\end{columns}
\end{frame}





\begin{frame}[fragile]
\frametitle{Símbolos Griegos}

\begin{columns}
    \begin{column}{0.5\textwidth}
        \begin{minted}{latex}
\[\sin \phi\]
        \end{minted}
    \end{column}
    {\color{gray} \vline}
    \begin{column}{0.5\textwidth}
        \[\sin \phi\]
    \end{column}
\end{columns}
\end{frame}





\begin{frame}[fragile]
\frametitle{Símbolos Griegos}

\begin{columns}
    \begin{column}{0.5\textwidth}
        \begin{tabular}{cl}
            \toprule
            \(\alpha A\) & \mintinline{latex}{\alpha A} \\
            \(\beta B\) & \mintinline{latex}{\beta B} \\
            \(\gamma \Gamma\) & \mintinline{latex}{\gamma \Gamma} \\
            \(\delta \Delta\) & \mintinline{latex}{\delta \Delta} \\
            \(\epsilon \varepsilon E\) & \mintinline{latex}{\epsilon \varepsilon E} \\
            \(\zeta Z\) & \mintinline{latex}{\zeta Z} \\
            \(\eta H\) & \mintinline{latex}{\eta H} \\
            \(\theta \vartheta \Theta\) & \mintinline{latex}{\theta \vartheta \Theta} \\
            \(\iota I\) & \mintinline{latex}{\iota I} \\
            \(\kappa K\) & \mintinline{latex}{\kappa K} \\
            \(\lambda \Lambda\) & \mintinline{latex}{\lambda \Lambda} \\
            \(\mu M\) & \mintinline{latex}{\mu M} \\
            \bottomrule
        \end{tabular}
    \end{column}
    \begin{column}{0.5\textwidth}
        \begin{tabular}{cl}
            \toprule
             \(\nu N\) & \mintinline{latex}{\nu N} \\
             \(\xi \Xi\) & \mintinline{latex}{\xi \Xi} \\
             \(o O\) & \mintinline{latex}{o O} \\
             \(\pi \Pi\) & \mintinline{latex}{\pi \Pi} \\
             \(\rho \varrho P\) & \mintinline{latex}{\rho \varrho P} \\
             \(\sigma \Sigma\) & \mintinline{latex}{\sigma \Sigma} \\
             \(\tau T\) & \mintinline{latex}{\tau T} \\
             \(\upsilon \Upsilon\) & \mintinline{latex}{\upsilon \Upsilon} \\
             \(\phi \varphi \Phi\) & \mintinline{latex}{\phi \varphi \Phi} \\
             \(\chi X\) & \mintinline{latex}{\chi X} \\
             \(\psi \Psi\) & \mintinline{latex}{\psi \Psi} \\
             \(\omega \Omega\) & \mintinline{latex}{\omega \Omega} \\
            \bottomrule
        \end{tabular}
    \end{column}
\end{columns}
\end{frame}





\begin{frame}[fragile]
\frametitle{Fuentes matemáticas}

\begingroup
\centering
\begin{tabular}{lcl}
    \toprule
    Tipo & Visualización & Comando \\
    \midrule
    Normal & \(\mathnormal{ABC}\) & \mintinline{latex}{\mathnormal{ABC}} \\
    Roman & \(\mathrm{ABC}\) & \mintinline{latex}{\mathrm{ABC}} \\
    Bold & \(\mathbf{ABC}\) & \mintinline{latex}{\mathbf{ABC}} \\
    Calligraphic & \(\mathcal{ABC}\) & \mintinline{latex}{\mathcal{ABC}} \\
    Sans Serif & \(\mathsf{ABC}\) & \mintinline{latex}{\mathsf{ABC}} \\
    Italic & \(\mathit{ABC}\) & \mintinline{latex}{\mathit{ABC}} \\
    Typewritter & \(\mathtt{ABC}\) & \mintinline{latex}{\mathtt{ABC}} \\
    \bottomrule
\end{tabular}
\par
\endgroup
\end{frame}





\section{El estándar de la \hologo{AmS}}

\begin{frame}
\frametitle{¿Por qué no basta con el modo básico?}
La \textit{American Mathematical Society} (\hologo{AmS}) desarrolló una serie de paquetes para resolver problemas que {\lmr\TeX} no manejaba de forma nativa:

\bigskip\pause

\begin{itemize}[<+->]
    \item \textbf{Alineación:} Ecuaciones que ocupan varias líneas.
    \item \textbf{Numeración:} Control fino sobre qué se numera y cómo.
    \item \textbf{Símbolos:} Hay símbolos matemáticos que no existen (\(\mathbb{R}\)).
    \item \textbf{Estructura:} Entornos específicos para teoremas, lemas y pruebas.
\end{itemize}
\end{frame}





\begin{frame}[fragile]
\frametitle{¿Por qué no basta con el modo básico?}
\begin{minted}{latex}
\usepackage{amsmath}
\usepackage{amssymb}
\end{minted}
\end{frame}





\begin{frame}[fragile]
\frametitle{¿Por qué no basta con el modo básico?}

Cuando el paquete \texttt{amsmath} es incluido, la definición de \mintinline{latex}{\[\]} cambia para ser un alias del entorno \texttt{equation*}.

El entorno \texttt{equation*} es similar a \texttt{displaymath}, pero es mucho más personalizable.
\end{frame}





\section{Alineación de Ecuaciones}


\subsection{\texttt{align}}

\begin{frame}[fragile]
\frametitle{El entorno \texttt{align}}

Permite alinear en columnas usando el símbolo \emph{et} (\texttt{\&}).

\begin{columns}
    \begin{column}{0.5\textwidth}
        \begin{minted}{latex}
\begin{align}
  (a+b)^2 &= a^2 + 2ab + b^2 \\
  (a-b)^2 &= a^2 - 2ab + b^2
\end{align}
        \end{minted}
    \end{column}
    {\color{gray} \vline}
    \begin{column}{0.5\textwidth}
        \lmr
        \begin{displaymath}
            \begin{aligned}
                (a+b)^2 &= a^2 + 2ab + b^2 \quad \text{(1)} \\
                (a-b)^2 &= a^2 - 2ab + b^2 \quad \text{(2)}
            \end{aligned}
        \end{displaymath}
    \end{column}
\end{columns}

\bigskip

\small \textit{Nota: El símbolo \texttt{\&} define el punto de anclaje (usualmente el igual).}
\end{frame}





\begin{frame}[fragile]
\frametitle{El entorno \texttt{align}}
Permite alinear en columnas usando el símbolo \emph{et} (\texttt{\&}).

\begin{columns}
    \begin{column}{0.5\textwidth}
        \begin{minted}{latex}
\begin{align*}
  \int \frac{2x}{x^2} \, dx
    &= \int \frac{1}{u} \, du
    & u &= x^2 \\
  &= \ln(u) + C \\
  &= \ln(x^2) + C
\end{align*}
        \end{minted}
    \end{column}
    {\color{gray} \vline}
    \begin{column}{0.5\textwidth}
        \begin{displaymath}
            \begin{aligned}
                \int \frac{2x}{x^2} \, dx
                    &= \int \frac{1}{u} \, du
                    & u &= x^2 \\
                &= \ln(u) + C \\
                &= \ln(x^2) + C
            \end{aligned}
        \end{displaymath}
    \end{column}
\end{columns}
\end{frame}





\subsection{\texttt{gather}}

\begin{frame}[fragile]
\frametitle{El entorno \texttt{gather}}
Centra varias ecuaciones sin alinearlas entre sí.

\begin{columns}
    \begin{column}{0.5\textwidth}
        \begin{minted}{latex}
\begin{gather}
  a^2 \pm 2ab + b^2
    = (a \pm b)^2 \\
  a^2 - b^2 = (a + b)(a - b)
\end{gather}
        \end{minted}
    \end{column}
    {\color{gray} \vline}
    \begin{column}{0.5\textwidth}
        \lmr
        \begin{displaymath}
            \begin{gathered}
                a^2 \pm 2ab + b^2 = (a \pm b)^2 \quad \text{(1)} \\
                a^2 - b^2 = (a + b)(a - b) \quad \text{(2)}
            \end{gathered}
        \end{displaymath}
    \end{column}
\end{columns}
\end{frame}





\begin{frame}[fragile]
\frametitle{El entorno \texttt{gather}}
Centra varias ecuaciones sin alinearlas entre sí.

\begin{columns}
    \begin{column}{0.5\textwidth}
        \begin{minted}{latex}
\begin{gather*}
  \int \frac{2x}{x^2} \, dx
    = \int \frac{1}{u} \, du \\
  u = x^2 \\
  = \ln(u) + C \\
  = \ln(x^2) + C
\end{gather*}
        \end{minted}
    \end{column}
    {\color{gray} \vline}
    \begin{column}{0.5\textwidth}
        \begin{displaymath}
            \begin{gathered}
                \int \frac{2x}{x^2} \, dx
                    = \int \frac{1}{u} \, du \\
                u = x^2 \\
                = \ln(u) + C \\
                = \ln(x^2) + C
            \end{gathered}
        \end{displaymath}
    \end{column}
\end{columns}
\end{frame}





\subsection{\texttt{multline}}

\begin{frame}[fragile]
\frametitle{El entorno \texttt{multline}}
Para una sola ecuación que es demasiado larga para una línea.

\begin{columns}
    \begin{column}{0.5\textwidth}
        \begin{minted}{latex}
\begin{multline}
  A = a + b + c + d + e \\
  + f + g + h + i + j + k \\
  + l + m + n + o + p
\end{multline}
        \end{minted}
    \end{column}
    {\color{gray} \vline}
    \begin{column}{0.5\textwidth}
        \lmr
        \begin{multline*}
            A = a + b + c + d + e \\
            + f + g + h + i + j + k \\
            + l + m + n + o + p \quad \text{(1)}
        \end{multline*}
    \end{column}
\end{columns}
\end{frame}





\begin{frame}[fragile]
\frametitle{El entorno \texttt{multline}}
Para una sola ecuación que es demasiado larga para una línea.

\begin{columns}
    \begin{column}{0.5\textwidth}
        \begin{minted}{latex}
\begin{multline*}
  A = a + b + c + d + e \\
  + f + g + h + i + j + k \\
  + l + m + n + o + p
\end{multline*}
        \end{minted}
    \end{column}
    {\color{gray} \vline}
    \begin{column}{0.5\textwidth}
        \begin{multline*}
            A = a + b + c + d + e \\
            + f + g + h + i + j + k \\
            + l + m + n + o + p
        \end{multline*}
    \end{column}
\end{columns}
\end{frame}





\subsection{\texttt{aligned}, \texttt{gathered}, \texttt{split}}

\begin{frame}[fragile]
\frametitle{Ecuaciones numeradas como bloque}

A veces queremos una sola numeración para varias líneas de desarrollo.

La \hologo{AmS} nos ofrece los entornos:

\begin{itemize}
    \item \texttt{aligned} \uncover<2->{\(\rightarrow\) \texttt{align}}
    \item \texttt{gathered} \uncover<2->{\(\rightarrow\) \texttt{gather}}
    \item \texttt{split} \uncover<2->{\(\rightarrow\) \texttt{align}\alert{*}}
\end{itemize}

\pause

Que funcionan igual que su contraparte, pero solamente se pueden utilizar dentro de otro entorno.
\end{frame}

\begin{frame}[fragile]
\frametitle{Ecuaciones numeradas como bloque}

\begin{columns}
    \begin{column}{0.5\textwidth}
        \begin{minted}{latex}
\begin{equation}
  \begin{aligned}
    a + b &= 7 \\
    ab &= 10
  \end{aligned}
\end{equation}
        \end{minted}
    \end{column}

    \begin{column}{0.5\textwidth}
        \begin{displaymath}
            \begin{aligned}
                a + b &= 7 \\
                ab &= 10
            \end{aligned}
            \quad (1)
        \end{displaymath}
    \end{column}
\end{columns}
\end{frame}





\section{Matrices y Arreglos}

\begin{frame}
\frametitle{Matrices de la \hologo{AmS}}

Los paquetes de la \hologo{AmS} facilitan la creación de matrices con diferentes delimitadores:
\begin{itemize}
    \item \texttt{matrix}: sin nada
    \item \texttt{pmatrix}: paréntesis
    \item \texttt{bmatrix}: corchetes
    \item ...
\end{itemize}
\end{frame}





\begin{frame}[fragile]
\frametitle{Matrices de la \hologo{AmS}}

\begin{columns}
    \begin{column}{0.5\textwidth}
        \begin{minted}{latex}
\begin{pmatrix}
  1 & 0 & 0 \\
  0 & 1 & 0 \\
  0 & 0 & 1
\end{pmatrix}
        \end{minted}
    \end{column}

    \begin{column}{0.5\textwidth}
        \begin{displaymath}
            \begin{pmatrix}
                1 & 0 & 0 \\
                0 & 1 & 0 \\
                0 & 0 & 1
            \end{pmatrix}
        \end{displaymath}
    \end{column}
\end{columns}
\end{frame}





\begin{frame}[fragile]
\frametitle{Matrices de la \hologo{AmS}}

\begin{columns}
    \begin{column}{0.5\textwidth}
        \begin{minted}{latex}
\begin{bmatrix}
  1 & 0 \\
  0 & 1 \\
  0 & 0
\end{bmatrix}
        \end{minted}
    \end{column}
    {\color{gray} \vline}
    \begin{column}{0.5\textwidth}
        \begin{displaymath}
            \begin{bmatrix}
                1 & 0 \\
                0 & 1 \\
                0 & 0
            \end{bmatrix}
        \end{displaymath}
    \end{column}
\end{columns}
\end{frame}





\subsection{Fuentes}

\begin{frame}[fragile]
\frametitle{Fuentes de la \hologo{AmS}}

\begingroup
\centering
\begin{tabular}{lcl}
    \toprule
    Tipo & Visualización & Comando \\
    \midrule
    Blackboard Bold & \(\mathbb{ABC}\) & \mintinline{latex}{\mathbb{ABC}} \\
    Gothic & \(\mathfrak{ABC}\) & \mintinline{latex}{\mathfrak{ABC}} \\
    \bottomrule
\end{tabular}
\par
\endgroup
\end{frame}





\section{Definiciones por casos}
\begin{frame}[fragile]
\frametitle{Definiciones por casos}
Fundamental para funciones definidas a trozos.

\begin{columns}
    \begin{column}{0.5\textwidth}
        \begin{minted}{latex}
f(x) = \begin{cases}
  1 & \text{si } x \in \mathbb{Q} \\
  0 & \text{si \(x \notin \mathbb{Q}\)}
\end{cases}
        \end{minted}
    \end{column}

    \begin{column}{0.5\textwidth}
        \begin{displaymath}
            f(x) = \begin{cases}
                1 & \text{si } x \in \mathbb{Q} \\
                0 & \text{si  \(x \notin \mathbb{Q}\)}
            \end{cases}
        \end{displaymath}
    \end{column}
\end{columns}
\end{frame}





\begin{frame}
\frametitle{Trucos y Datos Curiosos}

Existen calculadoras y programas que son capaces de recibir o recibir matemáticas escritas en el formato de {\lmr\LaTeX}.

\begin{itemize}[<+->]
    \item \alert<1>{\bf Macsyma} o \alert<1>{\bf Maxima}: Es un CAS muy potente, que es capaz de entregar las operaciones en el formato de {\lmr\LaTeX}.
    \item \alert<2-3>{\bf Emacs}, un editor de código de GNU con el que podemos editar un documento de {\lmr\LaTeX} (\emph{\lmr AUC\TeX}). Además es capaz de leer y operar sobre ecuaciones escritas en el mismo, mediante \alert<3>{\bf Calc Mode}, que es un CAS.
\end{itemize}

\pause

Por lo que podemos realizar cálculos complejos en una calculadora (que entiende álgebra), y obtener la representación en {\lmr\LaTeX} de lo que estamos calculando.
\end{frame}





\begin{frame}[allowframebreaks]
\frametitle{Referencias}
\nocite{*}
\printbibliography
\end{frame}
\end{document}
